\chapter{无监督学习与降维}

\section{本章导读：没有标签时，算法只能提出结构候选}

无监督学习很有诱惑力，因为它似乎能让数据自己“说话”。但数据不会自己解释自己。PCA、聚类、密度方法和降维可视化能够提出结构候选，却不能自动证明这些结构对应真实物理类别、演化路径或新天体族群。

本章把无监督学习放在探索性分析的位置上：它适合发现可疑结构、压缩高维特征、给后续监督任务提供线索，也适合帮助我们看见预处理尺度和样本选择造成的形状。但它的输出必须回到原始特征、数据质量和物理背景中检查。

读这一章时，请避免给聚类结果过早命名。一个 cluster 首先只是算法在特征空间里看到的一片区域；它是否有物理意义，还需要图表、样本复核、领域知识和后续实验共同支持。

\section{学习目标}

学完本章后，你应该能够：

\begin{itemize}
    \item 区分监督学习与无监督学习各自在回答什么问题；
    \item 理解 PCA 为什么可以用更少的坐标描述主要结构，并知道它本质上是在寻找方差最大的方向；
    \item 了解 $k$-means、DBSCAN 和 Gaussian mixture 各自适合什么样的聚类情形；
    \item 理解 t-SNE 与 UMAP 更适合做可视化探索，而不是直接当作科学结论；
    \item 知道无监督学习在 HR 图结构分析、光谱分群和异常候选初筛中的价值与边界。
\end{itemize}

\section{训练目标与贯穿微型项目}

本章的微型项目是在 HR 图特征上探索无监督结构。正文负责解释 PCA、\(k\)-means、DBSCAN、Gaussian mixture、t-SNE/UMAP 以及“算法结构不等于自然类别”；训练材料则要求你把无标签探索写进 Model Experiment Record，并额外保留候选结构与人工复核入口。

\begin{examplebox}[最小无监督项目结构]
\begin{lstlisting}[style=terminal]
unsupervised_hr_demo/
  data/
    stellar_hr_unsupervised_demo.csv
  notebooks/
    ch23_unsupervised_learning.ipynb
  figures/
    unsupervised_structure_map.png
  results/
    model_experiment_record.md
    structure_candidates.csv
\end{lstlisting}
\end{examplebox}

最终交付应包含：标准化方式、降维结果、聚类候选、参数敏感性、与 HR 图物理结构的对应关系，以及不能把算法分组直接命名为自然类别的限制说明。

\section{无监督学习的两个主问题}

本章虽然会出现多种算法，但它们服务的主问题可以先压成两个：

\begin{enumerate}
    \item \textbf{降维}：哪些变化方向最大，哪些低维坐标能保留主要结构；
    \item \textbf{分群}：哪些样本在特征空间中彼此接近，哪些点远离主体结构。
\end{enumerate}

这两件事都不是物理命名。PCA、t-SNE/UMAP 和聚类算法能提出结构候选，但不能单独证明这些候选就是恒星演化阶段、天体族群或新物理类别。

\section{没有标签时，我们还能问什么}

前面的章节里，我们一直在做“给定输入，预测标签或数值”的任务。但在真实的天文与物理数据里，常常没有现成标签，或者标签本身很昂贵。这时我们仍然有很多值得问的问题：

\begin{itemize}
    \item 数据中是否自然分成若干群体；
    \item 哪些维度真正承载了主要变化；
    \item 有没有显著偏离主体结构的样本值得人工复核；
    \item 某些我们习惯用文字描述的“序列”或“支”，是否能被算法自动识别出来。
\end{itemize}

这类问题就是无监督学习的典型工作场景。它的目标不是“没有标签时勉强凑合”，而是在不预先规定答案的情况下，先让数据自己暴露结构。

\section{一个教学版 HR 图结构问题}

本章配套 notebook 使用 \nolinkurl{stellar_hr_unsupervised_demo.csv}。这个小数据集用颜色 \texttt{bp\_rp} 和绝对星等 \texttt{absolute\_g\_mag} 构造了一个教学版 HR 图，包含：

\begin{itemize}
    \item 主序星；
    \item 红巨星；
    \item 白矮星；
    \item 一个刻意放进去的异常亮蓝目标。
\end{itemize}

在这个例子里，我们不会把参考类别直接交给聚类算法，而是只在事后用它来解释聚类结果。这正符合无监督学习的工作方式：\textbf{算法先看结构，人再解释结构。}

\section{为什么“先标准化”在无监督学习里特别关键}

无监督学习比监督学习更容易受特征尺度影响，因为它通常直接依赖：

\begin{itemize}
    \item 方差大小；
    \item 样本间距离；
    \item 局部密度；
    \item 相似度结构。
\end{itemize}

如果一个维度量纲远大于其他维度，那么 PCA 会优先把它当作主方向，$k$-means 会优先沿它分簇，DBSCAN 的邻域判定也会受到它的主导。因此，本章 notebook 先对特征做标准化，再进入 PCA 和聚类。

\section{PCA：找出变化最大的方向}

很多高维数据里，真正重要的变化方向并没有你想象得那么多。\textbf{主成分分析（PCA）}做的事情，就是在保持主要变化趋势的同时，寻找一组新的坐标轴，把数据投影到更低维空间。

若中心化后的数据协方差矩阵记为 $\Sigma$，PCA 要找的是协方差矩阵的主特征向量。对单位向量 $w$，投影方差可以写成

\[
\mathrm{Var}(w^\top x) = w^\top \Sigma w.
\]

第一主成分就是让这个量最大的方向：

\[
w_1 = \arg\max_{\|w\|=1} w^\top \Sigma w.
\]

它的物理直觉非常朴素：\textbf{先找出数据变化最剧烈的方向。}

\begin{examplebox}[PCA 的核心想法]
\begin{lstlisting}[style=pythonstyle]
original features --> rotate axes --> keep the directions
that explain the largest variance
\end{lstlisting}
\end{examplebox}

\section{PCA 的最小计算流程}

实际使用 PCA 时，可以把流程拆成四步：

\begin{enumerate}
    \item 对每个特征做中心化或标准化；
    \item 计算特征之间的协方差矩阵；
    \item 求协方差矩阵的特征向量和特征值；
    \item 按特征值从大到小选择主成分，并把样本投影到这些方向上。
\end{enumerate}

\begin{examplebox}[PCA 流程骨架]
\begin{lstlisting}[style=pythonstyle]
scaled_features = standardize(features)
covariance = compute_covariance_matrix(scaled_features)
eigenvalues, eigenvectors = eigendecompose(covariance)
components = choose_largest_eigenvectors(eigenvalues, eigenvectors)
projected = project(scaled_features, components)
\end{lstlisting}
\end{examplebox}

这段骨架的教学重点不是让你手写所有线性代数函数，而是让你明白 PCA 不是“自动降维按钮”。它先定义尺度，再寻找方差方向，最后保留一部分方向。任何一步改变，都会改变最终二维图或后续聚类结果。

\section{PCA 在 HR 图里意味着什么}

在教学版 HR 图中，颜色和绝对星等并不是独立变化的。主序星、巨星和白矮星在二维平面里形成有方向性的结构。因此，PCA 能帮助我们：

\begin{itemize}
    \item 识别“主要变化轴”；
    \item 用更少维度表达大部分结构；
    \item 为后续聚类和可视化提供更紧凑的表示。
\end{itemize}

但必须注意：PCA 找到的是\textbf{统计方差最大方向}，不一定自动等于某个最具物理意义的变量方向。统计结构和物理解释之间，仍然需要研究者自己搭桥。

\section{本章算法卡片的最小问题}

无监督学习很容易变成方法名列表，所以每个算法都应至少回答一个“它相信什么结构”的问题：

\begin{itemize}
    \item \textbf{PCA}：相信最大方差方向能概括主要变化；要检查中心化、协方差、解释方差和被丢弃的信息；
    \item \textbf{\(k\)-means}：相信样本可以围绕若干中心形成近似球状簇；要检查 \(k\)、初始化和簇内平方距离；
    \item \textbf{DBSCAN}：相信密度连通结构有意义；要检查 \texttt{eps}、\texttt{min\_samples}、噪声点和尺度敏感性；
    \item \textbf{Gaussian mixture}：相信数据可由若干概率成分混合而成；要检查软归属、协方差结构和边界样本；
    \item \textbf{t-SNE/UMAP}：主要用于可视化邻近关系；不能把二维图上的间距或团块直接当成科学结构。
\end{itemize}

这些问题可以写进 Algorithm Card，也可以作为本章 Model Experiment Record 的 chapter-specific fields。关键不是多填一张表，而是让读者知道算法到底根据什么把样本放在一起。

\begin{examplebox}[Algorithm Card: PCA]
\begin{itemize}
    \item \textbf{任务类型}：降维、结构探索、可视化前处理；
    \item \textbf{模型形式}：把原特征旋转到一组正交主成分方向；
    \item \textbf{优化目标}：第一主成分最大化投影方差，后续主成分在正交约束下继续最大化剩余方差；
    \item \textbf{用户固定}：标准化方式、保留主成分数、是否先做质量筛选；
    \item \textbf{诊断}：解释方差比例、载荷向量、投影图和被丢弃维度的信息；
    \item \textbf{失败模式}：尺度主导、非线性结构被压平、低方差信号被误删；
    \item \textbf{科学边界}：能说明主要统计变化方向，不能自动说明最重要物理机制。
\end{itemize}
\end{examplebox}

\begin{examplebox}[Algorithm Card: \texorpdfstring{$k$-means}{k-means}]
\begin{itemize}
    \item \textbf{任务类型}：无监督分群；
    \item \textbf{模型形式}：用 \(k\) 个中心代表 \(k\) 个簇，每个样本归到最近中心；
    \item \textbf{优化目标}：最小化簇内平方距离和；
    \item \textbf{用户固定}：簇数 \(k\)、初始化方式、距离度量和特征缩放；
    \item \textbf{诊断}：簇内距离、不同初始化结果、边界样本和每簇样本数；
    \item \textbf{失败模式}：强制所有点入簇、偏好近似球状簇、对尺度和初始化敏感；
    \item \textbf{科学边界}：能给出快速结构候选，不能单独证明簇是物理类别。
\end{itemize}
\end{examplebox}

\begin{examplebox}[Algorithm Card: DBSCAN]
\begin{itemize}
    \item \textbf{任务类型}：密度聚类、噪声点识别；
    \item \textbf{决策规则}：若点的 \texttt{eps} 邻域内至少有 \texttt{min\_samples} 个样本，则可成为密度连通结构的一部分；
    \item \textbf{用户固定}：\texttt{eps}、\texttt{min\_samples}、距离度量和标准化方式；
    \item \textbf{诊断}：簇数量、噪声点比例、参数敏感性和边界点位置；
    \item \textbf{失败模式}：不同密度结构难以同时处理，\texttt{eps} 选择会强烈改变结果；
    \item \textbf{科学边界}：能标出密度结构和候选噪声点，但噪声点不等于新天体。
\end{itemize}
\end{examplebox}

\begin{examplebox}[Algorithm Card: Gaussian Mixture Model]
\begin{itemize}
    \item \textbf{任务类型}：概率聚类、软归属建模；
    \item \textbf{模型形式}：用多个高斯成分的加权和近似数据分布；
    \item \textbf{优化目标}：估计每个成分的权重、均值和协方差，使观测数据似然较高；
    \item \textbf{用户固定}：成分数、协方差形式、初始化和收敛设置；
    \item \textbf{诊断}：样本软归属概率、边界样本、协方差椭圆和成分稳定性；
    \item \textbf{失败模式}：非高斯结构被硬套成高斯成分，成分数选择不稳；
    \item \textbf{科学边界}：能表达模糊边界，不保证每个概率成分对应真实物理族群。
\end{itemize}
\end{examplebox}

\begin{examplebox}[Algorithm Card: t-SNE / UMAP]
\begin{itemize}
    \item \textbf{任务类型}：高维数据二维可视化；
    \item \textbf{决策规则}：尽量在低维图上保留原空间中的局部邻近关系；
    \item \textbf{用户固定}：邻域大小、随机种子、距离度量和其他可视化超参数；
    \item \textbf{诊断}：不同随机种子和参数下结构是否稳定，是否回到原始特征验证；
    \item \textbf{失败模式}：二维团块被过度解释，团块间距离被误读为物理距离；
    \item \textbf{科学边界}：适合提出可视化假设，不适合单独作为定量分类证据。
\end{itemize}
\end{examplebox}

\section{把无监督算法先看成三类假设}

进入具体工具前，可以先把本章算法分成三类假设。

\begin{itemize}
    \item \textbf{方差假设}：PCA 相信大方差方向保留了主要结构；
    \item \textbf{距离/中心假设}：\(k\)-means 相信样本围绕若干中心组织；
    \item \textbf{密度/概率假设}：DBSCAN 相信密度连通有意义，GMM 相信样本来自若干概率成分。
\end{itemize}

这三类假设没有谁天然更高级。HR 图里的主序结构可能更像连续带状结构，星系颜色空间可能出现重叠概率成分，光谱特征里也可能同时存在中心群体和孤立异常点。选择算法时，首先要问：我愿意让算法相信哪一种结构？

\paragraph{解释方差不是解释物理}

PCA 常报告 explained variance ratio，也就是每个主成分解释了总方差的多少。例如前两个主成分解释了 \(90\%\) 方差，说明二维投影保留了大部分数值变化。但这不等于“保留了 \(90\%\) 物理意义”。如果一个罕见但重要的物理信号只占很小方差，它仍可能在 PCA 压缩中被削弱。因此 PCA 报告里应同时写：保留了多少方差，以及哪些科学上重要的低方差信息可能被牺牲。

\section{\texorpdfstring{$k$-means}{k-means}：把样本分成若干最近中心群}

\textbf{$k$-means} 的直觉很简单：先假设数据要分成 $k$ 群，再不断重复两步：

\begin{enumerate}
    \item 把每个样本分给最近的中心；
    \item 用各组样本的均值更新中心。
\end{enumerate}

它本质上在最小化簇内平方和：

\[
\sum_{i=1}^{N}
\left\|x_i - \mu_{c_i}\right\|^2,
\]

其中 $\mu_{c_i}$ 是样本 $x_i$ 所属簇的中心。

它的优点是：

\begin{itemize}
    \item 实现简单；
    \item 速度快；
    \item 对球状、大小接近的簇常常很有效。
\end{itemize}

它的限制也非常明确：

\begin{itemize}
    \item 必须提前指定 $k$；
    \item 对初始化敏感；
    \item 会强制每个点都属于某个簇；
    \item 对细长结构、不同密度结构和孤立异常点不够自然。
\end{itemize}

\section{\texorpdfstring{$k$-means}{k-means} 的一次迭代长什么样}

为了让算法不只停留在语言描述上，可以把一次 $k$-means 迭代写成：

\begin{examplebox}[\texorpdfstring{$k$-means}{k-means} 迭代骨架]
\begin{lstlisting}[style=pythonstyle]
centers = initialize_centers(rows, k)
for step in range(max_steps):
    assignments = assign_each_row_to_nearest_center(rows, centers)
    new_centers = recompute_centers(rows, assignments)
    if new_centers == centers:
        break
    centers = new_centers
\end{lstlisting}
\end{examplebox}

这段伪代码揭示了两个常被忽略的事实。第一，初始化会影响结果；第二，算法总会给每个样本一个簇标签，即使某个样本其实更像离群点。后者正是为什么异常候选或稀疏边界样本不宜直接交给 $k$-means 做最终判断。

\paragraph{为什么 \texorpdfstring{$k$-means}{k-means} 偏好球状簇}

\(k\)-means 的目标函数使用的是到中心的平方距离。这意味着一个簇被默认表示成“围绕某个中心的一团点”。如果真实结构是细长弧线、分叉结构或不同密度的连续带，中心加平方距离就会强行把它切成几块。天文数据里这很常见：HR 图主序不是一个圆团，光谱序列也未必围绕几个中心。因此 \(k\)-means 适合做快速粗分群，但不适合直接给复杂物理结构命名。

\section{DBSCAN：按密度而不是按中心来找结构}

\textbf{DBSCAN} 的思路和 $k$-means 很不一样。它不问“最近中心是谁”，而是问“附近是否有足够多的邻居”。因此它有两个非常有用的特点：

\begin{itemize}
    \item 可以把稀疏孤立点标成噪声；
    \item 不必事先指定簇数。
\end{itemize}

这使它更适合做：

\begin{itemize}
    \item 异常候选初筛；
    \item 非球状结构识别；
    \item 带有明显稀疏背景的数据探索。
\end{itemize}

在本章 HR 图教学数据里，DBSCAN 能比较自然地把主序星、红巨星和白矮星分开，并把那个人工放进去的异常亮蓝目标标记成噪声候选。这正好说明：\textbf{密度型聚类更愿意承认“这点不属于主体结构”。}

\section{DBSCAN 参数不是装饰}

DBSCAN 最重要的两个参数通常是邻域半径 \texttt{eps} 和最小邻居数 \texttt{min\_samples}。它们决定了什么叫“足够密集”：

\begin{itemize}
    \item \texttt{eps} 太小，许多正常点也会被标成噪声；
    \item \texttt{eps} 太大，不同结构可能被连成一团；
    \item \texttt{min\_samples} 太小，随机波动也可能形成簇；
    \item \texttt{min\_samples} 太大，小而真实的群体可能被忽略。
\end{itemize}

因此，无监督学习报告不能只写“用了 DBSCAN”。至少要记录参数、标准化方式、簇数量、噪声点数量，以及这些结果对参数变化是否敏感。

\paragraph{core、border 和 noise}

DBSCAN 的结果可以先分成三种点：

\begin{itemize}
    \item \textbf{core point}：邻域里样本足够多，是密度结构的骨架；
    \item \textbf{border point}：自己邻居不够多，但靠近某个 core point；
    \item \textbf{noise point}：既不是 core，也不贴近任何密度结构。
\end{itemize}

这个划分比“有几个簇”更值得学生记住。很多科学候选恰恰出现在 border 或 noise 区域，但它们可能是新奇目标，也可能只是测量质量差。因此 DBSCAN 给出的 noise label 应进入复核清单，而不是直接进入发现结论。

\section{Gaussian mixture：允许软归属}

$k$-means 默认每个样本只能属于一个簇，而且每个簇都像是“离哪个中心最近”这样的硬划分。\textbf{Gaussian mixture model（GMM）} 则更灵活：

\begin{itemize}
    \item 每个簇可以有不同形状和方向；
    \item 一个样本可以同时对多个簇有不同概率的归属；
    \item 更适合处理边界模糊、簇形状椭圆化的数据。
\end{itemize}

因此，如果你面对的是多种天体类型彼此过渡、测量误差较大、边界并不锐利的场景，GMM 往往比硬聚类更贴近实际。

GMM 的输出不是只有一个簇标签，而是一组归属概率。例如某个对象可能对成分 A 的概率为 \(0.55\)，对成分 B 的概率为 \(0.40\)。这类样本比概率 \(0.98\) 的样本更像边界对象，应该优先进入解释和复核。对科学分类来说，软归属的价值正在这里：它承认边界本身可能模糊。

\section{t-SNE 与 UMAP 更像“看图工具”}

t-SNE 和 UMAP 在现代天文数据分析里非常常见，因为它们常常能把高维数据压到二维，并把局部邻近关系展示得很清楚。但这里有一个特别重要的提醒：

\begin{itemize}
    \item 它们主要是\textbf{可视化工具}；
    \item 图上簇与簇之间的距离不一定有直接可解释的物理含义；
    \item 不同超参数会显著改变画面形状；
    \item “看起来分得很开”不等于科学上真的存在离散类别。
\end{itemize}

所以，t-SNE/UMAP 很适合辅助你发现结构、提出假设，却不适合单独承担定量结论。

\section{为什么无监督学习特别适合做初筛}

在天文与物理项目里，我们常常并不是想让无监督学习直接给出最终标签，而是希望它帮助我们：

\begin{itemize}
    \item 先把大样本整理成更容易检查的结构；
    \item 找出常见群体和少见群体；
    \item 把最值得人工复核的离群点排在前面。
\end{itemize}

这意味着无监督学习最适合做的是：

\begin{itemize}
    \item 结构探索；
    \item 样本整理；
    \item 新奇目标候选生成。
\end{itemize}

真正的物理解释、类别命名和后续验证，仍然必须回到观测背景、误差分析和领域知识。

\section{一个最小无监督结果报告模板}

无监督学习的报告不能写成“算法发现了三类天体”。更稳妥的写法是把算法结构和科学解释分开。一个最小报告应包含：

\begin{enumerate}
    \item \textbf{输入空间}：使用了哪些特征，是否标准化，是否剔除了明显质量问题；
    \item \textbf{方法与参数}：PCA 保留了几个主成分，聚类方法和关键参数是什么；
    \item \textbf{结构描述}：算法给出了多少簇、多少噪声点、哪些样本处在边界；
    \item \textbf{解释证据}：这些结构是否对应 HR 图位置、光谱形态、颜色范围或其他领域知识；
    \item \textbf{解释边界}：哪些结构只是候选假设，哪些需要后续标签、观测或人工复核。
\end{enumerate}

\begin{examplebox}[无监督报告的最小骨架]
\begin{lstlisting}[style=pythonstyle]
model_experiment_record["ch23_fields"] = {
    "features": ["bp_rp", "absolute_g_mag"],
    "scaling": "standardized before PCA and clustering",
    "method": "DBSCAN",
    "key_parameters": {"eps": eps, "min_samples": min_samples},
    "cluster_summary": cluster_counts,
    "noise_candidates": noise_ids,
    "interpretation_status": "candidate structure, not final labels",
}
\end{lstlisting}
\end{examplebox}

这种写法看起来谨慎，但并不削弱无监督学习的价值。相反，它让无监督结果真正成为可继续推进的科学假设：算法负责暴露结构，研究者负责验证结构是否有物理意义。

本章的 Pass 要求是：任何 cluster 命名前，必须先写明 ``algorithmic group, not physical class unless externally validated''。中文报告可以写成：“这是算法分组，不是物理类别；除非有外部证据验证，否则不能命名为真实天体族群。”

在 Part III Model Experiment Packet 里，无监督报告通常不会直接变成最终物理结论，而是进入 error analysis artifact、Trust Statement 和后续人工复核队列。它的价值是提出“值得进一步检查的结构”，而不是替研究者完成类别命名。

\section{无监督探索的安全顺序}

一个稳妥的无监督探索流程可以固定为：

\begin{enumerate}
    \item 先写明输入特征和质量筛选；
    \item 对尺度敏感方法先做标准化；
    \item 用 PCA 或二维图理解主要结构，但不急着命名；
    \item 对比至少两类聚类或离群方法；
    \item 检查参数敏感性；
    \item 把算法结构与物理解释分开写入报告。
\end{enumerate}

这条顺序的目的，是把“看见一张漂亮图”降级为“形成一个可复查的候选结构”。无监督学习的强项是提出问题；真正的科学命名需要额外证据。

\section{一个最常见的风险：把算法结构误当成自然分类}

无监督学习最大的诱惑在于：它会很快给你一个“看起来很有结构”的结果。但你必须保持警惕：

\begin{itemize}
    \item 聚类结果是否只是尺度选择造成的；
    \item PCA 压缩后丢掉了什么信息；
    \item DBSCAN 标出的“噪声”是新奇天体，还是测量误差；
    \item 你看到的是物理群体，还是观测选择效应；
    \item 一个漂亮的二维图，是否掩盖了原空间中的连续过渡。
\end{itemize}

\begin{warnbox}[无监督学习中的常见误区]
\begin{itemize}
    \item 没有标准化就直接聚类或做 PCA；
    \item 看到二维投影分开，就认定存在离散类别；
    \item 把异常点直接当作新发现，而不先排查数据质量；
    \item 用聚类标签替代科学解释；
    \item 反复调参直到得到“好看”的图，再把它当成证据。
\end{itemize}
\end{warnbox}

\section{AI 助手如何帮助你}

在这一章，AI 助手适合帮助你：

\begin{itemize}
    \item 检查 PCA、$k$-means 或 DBSCAN 的代码流程是否清楚；
    \item 帮你整理不同聚类参数下的结果摘要；
    \item 讨论某个离群点更像物理新奇目标还是数据质量问题；
    \item 协助你把“算法发现了结构”与“我们如何解释结构”分成两层论证。
\end{itemize}

但簇的命名、异常点的物理解释、以及可视化结果是否足以支持科学结论，仍然需要你自己作判断。

\section{练习}

\begin{exercisebox}[基础练习]
解释为什么无监督学习里“先标准化”往往比监督学习里更容易影响最终结构。
\end{exercisebox}

\begin{exercisebox}[进阶练习]
从 PCA 的目标
\[
\max_{\|w\|=1} w^\top \Sigma w
\]
出发，用语言解释为什么第一主成分可以理解为“变化最剧烈的方向”。
\end{exercisebox}

\begin{exercisebox}[开放练习]
为一个你感兴趣的数据集写一页无监督探索方案，至少包含：特征标准化方式、降维方法、候选聚类方法、异常点判别策略，以及你最担心的误解释风险。
\end{exercisebox}

\section{本章小结}

无监督学习不是在“没有标签时勉强凑合”，而是在探索数据结构时非常重要的一类方法。PCA 帮你压缩和观察主要变化；$k$-means 擅长快速分群；DBSCAN 更适合识别密度结构和噪声；GMM 提供更柔和的概率归属；t-SNE/UMAP 则更适合辅助可视化。真正可靠的工作方式，是把这些工具当作提出问题、整理样本和筛选候选的方法，而不是替代科学解释的自动答案。
